\documentclass[10pt,a4paper]{amsart}
\usepackage[margin=19mm]{geometry}
\usepackage{amsmath,amssymb,mathtools}
\usepackage[scr=rsfs,cal=euler]{mathalfa}
\usepackage{xfrac}
\usepackage[unicode,hidelinks,bookmarks=false]{hyperref}
\newcommand{\ie}{i.e.\ }
\newcommand{\cf}{cf.\ }
\newcommand{\resp}{respectively\ }
\newcommand{\Subsection}{\subsection}
\providecommand{\nobreakdash}{\nobreak\hbox{-}}
\setlength{\emergencystretch}{2em}
\multlinegap0pt
\usepackage{bm}
\usepackage{bbm}
\usepackage{dsfont}
\usepackage{xspace}
\usepackage{cancel}
\usepackage{mathtools}
\usepackage{amsthm}
\makeatletter
\@ifundefined{proposition}{\newtheorem{proposition}{Proposition}}{}
\makeatother
\def\a{\alpha}
\def\l{\lambda}
\def\p{\partial}
\title[An isoperimetric inequality for the second Robin eigenvalue ]{An isoperimetric inequality for the second Robin eigenvalue of the Weighted Laplacian}
\author{Yi Gao, Kui Wang, Anqiang Zhu}
\date{Source: arXiv:2602.19456v1 (CC BY 4.0)}
\begin{document}
\maketitle

\noindent\emph{Source and licence.} An isoperimetric inequality for the second Robin eigenvalue of the Weighted Laplacian. Version arXiv:2602.19456v1 (CC BY 4.0), distributed under the Creative Commons Attribution 4.0 International licence, as recorded in the arXiv licence metadata for this exact version and shown on the abstract page https://arxiv.org/abs/2602.19456v1. Full rights notice: licenses/arxiv-2602.19456-CC-BY-4.0.txt.\par\smallskip

\noindent\textit{Editorial scope.} Counted unit: the Proposition whose source label is \texttt{prop 2.3} in the article (source0.tex lines 388--398, source label \texttt{prop 2.3}), delivered together with the article's own complete original proof of it (source0.tex lines 399--495, one \texttt{proof} environment of 97 lines). No mathematics is written, repaired or re-derived: statement and proof are the article's own text line for line. Required context retained uncounted: eq 2.2 (lines 276--281); eq 2.4 (lines 316--321); eq 2.5 (lines 323--328); eq 2.7 (lines 331--333); prop 2.2 (lines 341--347). Every definition, display and label of those lines is retained unchanged. Omitted from this package and not redistributed: 1--275, 282--315, 322--322, 329--330, 334--340, 348--387, 496--831. Nothing inside a retained excerpt is omitted or altered: every retained body is delivered exactly as the article's lines are written, so no omission receipt of any kind accompanies this package. Citations: the retained excerpts carry no citation command at all, so no bibliography entry of the article is transcribed and no reference is redistributed. Reference adaptation: none was needed and none was made; every reference inside the delivered unit resolves inside it, so no unresolved reference or citation is produced. Printed slips kept literal: no printed word, symbol, hypothesis or number of the counted statement or of its proof was altered. Layout and class: the article's own class is \texttt{amsart} with packages including geometry, amsmath, amsfonts, amssymb, graphicx, color, hyperref, xcolor; the wrapper is the lane's pinned \texttt{amsart} editorial wrapper at 10pt on A4, which restates the theorem-like environments the retained text needs and adds the article's own macro definitions that the retained text needs (\texttt{\textbackslash a}, \texttt{\textbackslash l}, \texttt{\textbackslash p}), with extra wrapper packages (bm, bbm, dsfont, xspace, cancel, mathtools). The delivered numbering is the unit's own. Assets: no retained span contains a figure environment, a graphics-inclusion command, a PDF-page inclusion, a table or a TikZ picture; no image file is copied into this package, the package declares no asset dependency and no image file is redistributed with it. Licence: version arXiv:2602.19456v1 of this article is distributed under the Creative Commons Attribution 4.0 International licence, as recorded in the arXiv licence metadata for this exact version and shown on the abstract page https://arxiv.org/abs/2602.19456v1. A full rights notice is delivered at licenses/arxiv-2602.19456-CC-BY-4.0.txt. Characters: every retained excerpt is pure ASCII, so the delivered engine, XeLaTeX with its default Latin Modern text font, reports no missing character.
\begin{align}\label{eq 2.2}
\begin{cases}
    -\Delta u-\nabla h \cdot \nabla u=\l u \quad & \text{in $B(R)$,}\\
\frac{\p u}{\p \nu}+\a u=0 \quad & \text{on $\p B(R)$,}
\end{cases}
\end{align}

\begin{equation} \label{eq 2.4}
\begin{cases}
w'' + \left( \dfrac{m-1}{r} + h'(r) \right)w' +\l w=0 & \text{in } (0,R), \\
w'(0) =0, \qquad  w'(R)+\alpha w(R) = 0.
\end{cases}
\end{equation}

\begin{equation} \label{eq 2.5}
\begin{cases}
w'' + \left( \dfrac{m-1}{r} + h'(r) \right)w' + (\l-\frac {m-1} {r^2}) w = 0 & \text{in } (0,R), \\
w(0)=0, \qquad  w'(R)+\alpha w(R) = 0.
\end{cases}
\end{equation}

\begin{align}\label{eq 2.7}
   g'' + \left( \dfrac{m-1}{r} + h' \right)g' + (\mu_1-\frac {m-1} {r^2}) g = 0 
\end{align}

\begin{proposition}\label{prop 2.2}
Let $g(r)$ be the  positive  eigenfunction of \eqref{eq 2.5}. If $\a\le 0$, then
\begin{enumerate}
\item $g'(r)>0$ for $r\in (0, R)$.
\item Assume further that $\a\ge -\frac 2 R$, then  $g'(r)\ge  -\a g(r)$ for $r\in (0, R]$.
\end{enumerate}
\end{proposition}

\begin{proposition}\label{prop 2.3}
Suppose $\a\le 0$ and $h'(r)>0$. Then the second Robin eigenfunctions of \eqref{eq 2.2} are given by 
\begin{align*}
u_i(x)=g(r) \frac{x_i}{r},\quad  i=1,2, \cdots,m,
\end{align*}
where $g(r):[0, R]\rightarrow [0,\infty)$ solves  
\begin{equation}\label{eq 2.11}
g''(r)+(\frac {m-1} {r}+h'(r)) g'(r)+\left(\l_{2,\a}(B(R); \gamma_h)-\frac {m-1} {r^2}\right) g(r)=0
\end{equation}
with boundary condition $g(0)=0$ and $g'(R)=-\a g(R)$. Here $\l_{2, \a}(B(R);\a)$ denotes the second Robin eigenvalue of \eqref{eq 2.2}. 
\end{proposition}

\begin{proof}
  Since the second eigenvalue $\l_{2, \a}(B(R);\gamma_h)$ of \eqref{eq 2.2} is  either the first eigenvalue $\mu_1$ of \eqref{eq 2.5}
 or the second eigenvalue $\tau_2$ of \eqref{eq 2.4},  it suffices to show
\begin{align*}
    \mu_1<\tau_2.
\end{align*}
Recall that an eigenfuction corresponding to the second eigenvalue of \eqref{eq 2.4} must change sign in $(0,R)$. Without loss of generality, assume $f(r)$ is positive in $(0, a)$ for some $a<R$ and $f(a)=0$. Then  $f'(a)<0$, and the same argument as in part (1) of Proposition \ref{prop 2.2} shows $f'(r)<0$ for $r\in(0, a)$.
Let  $f(r)$ be an eigenfunction of \eqref{eq 2.4}, i.e.
\begin{align}\label{eq 2.6}
 f'' + \left( \dfrac{m-1}{r} + h' \right)f' + \tau_2 f = 0  
\end{align}
for $r\in (0, R)$, with $f'(0)=0$ and $f'(R)+\a f(R)=0$.
Differentiating  \eqref{eq 2.6} gives
\begin{align}\label{eq 2.13}
    f'''+(\frac{m-1}{r}+h')f''+(\tau_{2}-\frac{m-1}{r^{2}}+h'')f'=0.
\end{align}
Set $\phi(r)=e^{h(r)}f'(r)$,  then 
\begin{align*}
    \phi'(r)=e^{h(r)}f''+h'\phi, 
\end{align*}
and
\begin{align*}
    \phi''(r)&=e^{h(r)}f'''+e^{h(r)}h'f''+h''\phi+h'\phi'\\
    &=e^{h}(-(\frac{m-1}{r}+h'(r))f''-(\tau_2-\frac{m-1}{r^{2}}+h'')f')+e^{h(r)}h'f''+h''\phi+h'\phi'\\
    &=e^{h}f''(-(\frac{m-1}{r}+h'(r))+h')-(\tau_{2}-\frac{m-1}{r^{2}}+h'')\phi+h''\phi+h'\phi'\\
    &=(\phi'(r)-h'\phi(r))(-\frac{m-1}{r})-(\tau_{2}-\frac{m-1}{r^{2}})\phi+h'\phi'\\
&=(-\frac{m-1}{r}+h')\phi'-(\tau_{2}-\frac{m-1}{r^{2}}-\frac{m-1}{r}h')\phi.
\end{align*}
Thus, \eqref{eq 2.13} implies
\begin{align}\label{eq 2.14}
    \phi''+(\frac{m-1}{r}-h')\phi'+(\tau_{2}-\frac{m-1}{r^{2}}-\frac{m-1}{r}h')\phi=0.
\end{align}
From  \eqref{eq 2.7}, we have
\begin{align*}
    0&=g'' + \left( \dfrac{m-1}{r} + h' \right) g'+ (\mu_{1}-\frac{m-1}{r^{2}}) g\\
    &=g''+(\frac{m-1}{r}-h')g'+(\mu_{1}-\frac{m-1}{r^{2}}-\frac{m-1}{r}h') g+2h'g'+\frac{m-1}{r}h'g.
\end{align*}
and the facts $g(r)>0$,  $h'(r)\ge 0$ and $g'(r)>0$ for $r\in (0,a)$, we obtain
\begin{align}\label{eq 2.15}
    0\geq  g''+(\frac{m-1}{r}-h^{'})g'+(\mu_{1}-\frac{m-1}{r^{2}}-\frac{m-1}{r}h') g.
\end{align}
Because $f'(r)<0$ for $r\in (0,a)$, we have $\phi(r)<0$ for $r\in(0, a)$. Multiplying \eqref{eq 2.15} by $\phi$  yields
\begin{align}\label{eq 2.16}
    0\leq \phi g''+\phi g'(\frac{m-1}{r}-h')+(\mu_{1}-\frac{m-1}{r^{2}}-\frac{m-1}{r}h') \phi g.
\end{align}
Subtracting equation \eqref{eq 2.14} multiplied by  $g$
from \eqref{eq 2.16} gives
\begin{align}\label{eq 2.17}
    0\leq \phi g''-g\phi''+(\frac{m-1}{r}-h')(g'\phi-g\phi')+(\mu_{1}-\tau_{2})g\phi.
\end{align}
Multiplying \eqref{eq 2.17} by $r^{m-1}e^{-h}$ and integrating over $(0,a)$, we obtain
\begin{align}\label{eq 2.18}
    0\leq \int_{0}^{a}(\phi g''-g\phi''+(\frac{m-1}{r}-h')(g'\phi-g\phi')+(\mu_{1}-\tau_{2})g\phi)r^{m-1}e^{-h(r)}dr.
\end{align}
Direct computation yields
\begin{align*}
   & \int_{0}^{a}(\phi g''-g\phi'')r^{m-1}e^{-h(r)}dr\\
    =&\int_{0}^{a}(\phi g'r^{m-1}e^{-h(r)})'-\phi'g'r^{m-1}e^{-h}-(\frac{m-1}{r}-h')\phi g'r^{m-1}e^{-h}dr\\
    &-\int_{0}^{a}(g\phi'r^{m-1}e^{-h(r)})'dr+\int_{0}^{a}g'\phi'r^{m-1}e^{-h}dr+\int_{0}^{a}(\frac{m-1}{r}-h')g\phi^{'}dr,
\end{align*}
and therefore
\begin{align}\label{eq 2.19}
\begin{split}
       &\int_{0}^{a}(\phi g''-\phi''g)r^{m-1}e^{-h(r)}dr\\
    =&(\phi g'-g\phi')r^{m-1}e^{-h}\Big|_{0}^{a}+\int_{0}^{a}(\frac{m-1}{r}-h')(g\phi'-\phi g')r^{m-1}e^{h}dr. 
\end{split}
\end{align}
Substituting \eqref{eq 2.19} into  \eqref{eq 2.18} yields 
\begin{align}\label{eq 2.20}
    0\leq (\phi g'-g\phi')r^{m-1}e^{-h}\Big|_{0}^{a}+\int_{0}^{a}(\mu_{1}-\tau_{2})g\phi r^{m-1}e^{-h(r)}dr.
\end{align}
Since $g(0)=0$ and $\phi(0)=e^{h(0)}f'(0)=0$, we have 
\begin{align*}
    (\phi g'-g\phi')r^{m-1}e^{-h}\Big|_{0}^{a}=\phi(a)g'(a)-g(a)\phi'(a).
\end{align*}
Recall that  $g(a)>0$, $g' (a)>0$ and $\phi(a)=e^{h(a)}f'(a)<0$. From \eqref{eq 2.6}, we obtain  
\begin{align*}
    f''(a)+(\frac{m-1}{a}+h'(a))f'(a)+\tau_{2}f(a)=0
\end{align*}
Since $f(a)=0$ and $f'(a)\leq 0$, it follows that
\begin{align*}
    f''(a)+h'(a)f'(a)=-\frac{m-1}{a}f'(a)> 0,
\end{align*}
and consequently
\begin{align*}
    \phi'(a)=e^{h(a)}(f''(a)+h'(a)f'(a))> 0.
\end{align*}
Hence
\begin{align}\label{eq 2.21}
    (\phi g'-g\phi')r^{m-1}e^{-h}\Big|_{0}^{a}= (\phi(a) g'(a)-g(a)\phi'(a))a^{m-1}e^{-h(a)}< 0.
\end{align}
Substituting  \eqref{eq 2.21} into  \eqref{eq 2.20} gives
\begin{align*}
    0< \int_{0}^{a}(\mu_{1}-\tau_{2})g\phi r^{m-1}e^{-h(r)}dr
\end{align*}
Since $\phi(r)<0$ for $r\in (0,a)$, we conclude $ \mu_{1}< \tau_{2}$. This completes  the proof of  Proposition  \ref{prop 2.3}.
\end{proof}

\end{document}
